\chapter{Options asiatiques, lookbacks, cliquets et options à départ différé}\label{ch:dv:asians-lookbacks-cliquets-and-forward-starts}

Un trésorier d'entreprise demande un call d'un an sur l'indice, puis le même call sur la moyenne de douze
cours mensuels. La deuxième cotation est moins chère de 40~\%. Personne au téléphone ne sait où sont passés les 40~\%, et le trésorier
soupçonne une astuce. Il n'y en a aucune. La moyenne de douze cours fluctue beaucoup moins que le dernier~: sa volatilité vaut
environ 61~\% de celle de l'indice. Ce chapitre valorise les payoffs dépendant du chemin que les desks vendent le plus. Les asiatiques moyennent le
chemin, les lookbacks lisent ses extrêmes, et les \omterm{def:dv:asians-lookbacks-cliquets-and-forward-starts:fwdstart}{options à départ différé} et les \omterm{def:dv:asians-lookbacks-cliquets-and-forward-starts:cliquet}{cliquets} sont des options dont les strikes sont fixés dans le futur.
Les deux premières se valorisent bien avec tout modèle qui cale le smile d'aujourd'hui. Les deux dernières ne le sont pas, car leur valeur
dépend du smile que le modèle prédit pour les dates futures. Le chapitre mesure cette dépendance avec deux modèles qui
calent la même surface.

\section{Moyennes~: options asiatiques}

\begin{definition}[Option asiatique]\label{def:dv:asians-lookbacks-cliquets-and-forward-starts:asian}
Une \emph{option asiatique}\index{option asiatique} paie sur la moyenne du sous-jacent sur un ensemble de dates de fixation~: la
forme à moyenne de cours paie $(\bar S-K)^+$ (l'option sur prix moyen, livre~3, chapitre~12), et la forme à strike moyen
paie $(S_T-\bar S)^+$. La moyenne est arithmétique, sauf mention contraire.
\end{definition}

La moyenne réduit la volatilité. Avec $n$ fixations également espacées sur $T$, le logarithme de la moyenne géométrique suit une loi normale
de variance
\[
\sigma_G^2T=\sigma^2T\,\frac{(n+1)(2n+1)}{6n^2},
\]
qui tend vers $\sigma^2T/3$ lorsque les fixations deviennent continues. L'asiatique géométrique est donc une option de Black--Scholes
à volatilité et forward réduits, sous forme fermée. La moyenne arithmétique, celle qu'utilisent les contrats, n'a
pas de forme fermée. Elle est supérieure à la géométrique, mais les deux évoluent presque ensemble.

\begin{example}[Où sont passés les 40~\%]\label{ex:dv:asians-lookbacks-cliquets-and-forward-starts:asian}
Spot et strike 100, un an, $r=3\%$, $q=1\%$, volatilité 20~\%. Le call vanille vaut 8.83, et le call sur la
moyenne de douze fixations mensuelles vaut 5.32, soit 40~\% de moins. Le facteur de variance vaut $13\times25/(6\times144)=0.376$, donc la
moyenne géométrique fluctue avec une volatilité de $0.61\times20\%=12.3\%$. À mesure que les fixations se multiplient, le prix de l'asiatique rapporté
au call vanille baisse de 0.78 pour deux fixations à 0.68 pour quatre, 0.60 pour douze et 0.57 pour des fixations quotidiennes. Ce rapport suit
le facteur de volatilité $\sqrt{(n+1)(2n+1)/(6n^2)}$~: 0.79, 0.68, 0.61 et 0.58
(la \cref{fig:dv:asians-lookbacks-cliquets-and-forward-starts:asian}, à gauche).
\end{example}

La quasi-parfaite concordance des deux moyennes est ce qui rend peu coûteuse la valorisation par Monte-Carlo des asiatiques. L'asiatique géométrique
est une variable de contrôle (livre~4, chapitre~26). Kemna et Vorst l'ont proposée. Sur chaque trajectoire, calculer les deux payoffs, et
corriger l'estimation arithmétique par l'erreur connue du payoff géométrique~:
$\hat A_{\mathrm{cv}}=\bar A-\beta(\bar G-G_{\mathrm{exact}})$, avec $\beta$ estimé par régression sur les mêmes
trajectoires. Sur l'exemple, les deux payoffs actualisés ont une corrélation de 0.9996 et $\beta=1.03$. L'erreur type passe
de 0.025 à 0.0007, soit une réduction de variance d'un facteur 1\,357~: une trajectoire avec la variable de contrôle équivaut à 1\,357 trajectoires sans
elle (la \cref{fig:dv:asians-lookbacks-cliquets-and-forward-starts:asian}, à droite).

\begin{omfigure}
\begin{tikzpicture}
\begin{groupplot}[group style={group size=2 by 1,horizontal sep=1.6cm},
  width=6.6cm,height=5cm,
  tick label style={font=\small},label style={font=\small},grid=major,grid style={gray!25},table/col sep=comma]
\nextgroupplot[xmode=log,xlabel={fixations dans l'année},ylabel={rapport},xmin=0.8,xmax=320,ymin=0.5,ymax=1.05,
  xtick={1,2,4,12,52,252},xticklabels={1,2,4,12,52,252},
  title={\small Asiatique / vanille},
  legend columns=1,legend cell align=left,legend style={font=\footnotesize,at={(0.5,-0.32)},anchor=north,draw=none,fill=none}]
\addplot[omThm,very thick,mark=*,mark size=1.6pt] table[x=n,y=ratio]{figdata/derivatives/16-asians-lookbacks-cliquets-and-forward-starts/asian_fixings.csv};
\addplot[omDef,thick,dashed] table[x=n,y=volfactor]{figdata/derivatives/16-asians-lookbacks-cliquets-and-forward-starts/asian_fixings.csv};
\legend{rapport de prix,{$\sqrt{(n+1)(2n+1)/6n^2}$}}
\nextgroupplot[xlabel={payoff de la moyenne géométrique},ylabel={payoff de la moyenne arithmétique},xmin=0,xmax=40,ymin=0,ymax=40,
  title={\small 400 trajectoires, douze fixations}]
\addplot[omThm,only marks,mark=*,mark size=0.8pt] table[x=geometric,y=arithmetic]{figdata/derivatives/16-asians-lookbacks-cliquets-and-forward-starts/asian_scatter.csv};
\draw[gray,dashed] (axis cs:0,0) -- (axis cs:40,40);
\end{groupplot}
\end{tikzpicture}

\omcaption{À gauche~: le prix d'un call asiatique à la monnaie d'un an rapporté au call vanille, selon le nombre de fixations,
face au facteur de volatilité de la moyenne géométrique. À droite~: les payoffs arithmétique et géométrique, trajectoire par trajectoire~: presque
la même variable aléatoire, ce qui fait de l'asiatique géométrique, connue sous forme fermée, une variable de contrôle efficace. Données~:
le tutoriel.}\label{fig:dv:asians-lookbacks-cliquets-and-forward-starts:asian}
\end{omfigure}

Les asiatiques ont deux propriétés pratiques que le seul prix ne laisse pas voir. Leur risque diminue au fil des fixations~: une fois la moitié des
fixations sont constatées, la moitié du payoff est connue, et le delta et le \omterm{def:dv:greeks-and-the-hedging-pnl:greeks}{vega} se réduisent d'autant. La couverture au voisinage d'une fixation, en revanche,
concentre les flux à la date et à l'heure de la fixation, l'une des raisons pour lesquelles les moyennes portent sur de nombreuses fixations. Et la dynamique du smile
compte peu. La valeur d'une asiatique dépend surtout de la zone à la monnaie du smile à chaque date de fixation,
que tout modèle calé sur la surface reproduit. Sur la surface du chapitre~9 (taux nuls, douze fixations mensuelles),
l'asiatique à la monnaie vaut 4.34 sous \omterm{def:dv:local-volatility:model}{volatilité locale} et 4.36 sous Heston, à une erreur type près l'une de l'autre.
Black--Scholes à la volatilité à la monnaie d'un an donne 4.69~: ce qui compte est la structure par terme, puisque les premières
fixations voient des volatilités de courte échéance plus basses.

\section{Extrêmes~: lookbacks}

\begin{definition}[Option lookback]\label{def:dv:asians-lookbacks-cliquets-and-forward-starts:lookback}
Une \emph{option lookback}\index{option lookback} paie sur le maximum ou le minimum du sous-jacent sur sa durée de vie. Le
lookback call à strike flottant paie $S_T-\min_tS_t$ (achat au plus bas)~; le lookback call à strike fixe paie
$(\max_tS_t-K)^+$.
\end{definition}

Un lookback est une option à barrière intégrée sur tous les niveaux de barrière~: $S_T-\min S=\int\mathbf 1\{\min S<h\}\,dh$ plus
le prix terminal. Goldman, Sosin et Gatto l'ont valorisé sous forme fermée par le principe de réflexion, pour une surveillance
continue. Il est cher. Avec les paramètres de l'exemple, le lookback call à strike flottant vaut 15.69, soit 1.8 fois le
call à la monnaie, car il achète au meilleur cours de l'année.

Les contrats observent l'extrême à des dates discrètes, souvent les clôtures quotidiennes, et le minimum observé est supérieur au vrai
minimum. Le \omterm{def:dv:barriers-and-digitals:shift}{décalage de barrière} du chapitre~15 se transpose ici~: le minimum observé se comporte comme le minimum continu multiplié par
$e^{\beta\sigma\sqrt{\Delta t}}$, ce qui donne $C_{\mathrm d}\approx e^{x}C_{\mathrm c}-(e^x-1)Se^{-qT}$ avec
$x=\beta\sigma\sqrt{\Delta t}$. Avec un suivi quotidien, une simulation donne 15.09 (erreur type 0.04) et la formule décalée
15.08. Au rythme hebdomadaire, elles donnent 14.42 et 14.34~; mensuel, 13.13 et 12.84~; avec quatre dates, 11.55 et 10.70
(la \cref{fig:dv:asians-lookbacks-cliquets-and-forward-starts:lookback}). Comme pour les barrières, la correction s'applique au
suivi quotidien, et un calendrier clairsemé exige la simulation.

\begin{omfigure}
\begin{tikzpicture}
\begin{axis}[width=10.5cm,height=5cm,xmode=log,
  xlabel={dates de suivi par an},ylabel={lookback call à strike flottant},
  tick label style={font=\small},label style={font=\small},
  xmin=3,xmax=320,ymin=8,ymax=16.5,grid=major,grid style={gray!25},
  xtick={4,12,52,252},xticklabels={4,12,52,252},
  legend columns=4,legend style={font=\footnotesize,at={(0.5,-0.3)},anchor=north,draw=none,fill=none,/tikz/every even column/.append style={column sep=6pt}},
  table/col sep=comma]
\addplot[black,only marks,mark=*,mark size=1.8pt] table[x=n,y=mc]{figdata/derivatives/16-asians-lookbacks-cliquets-and-forward-starts/lookback.csv};
\addplot[omThm,very thick] table[x=n,y=shifted]{figdata/derivatives/16-asians-lookbacks-cliquets-and-forward-starts/lookback.csv};
\addplot[gray,thick,dashed] table[x=n,y=cont]{figdata/derivatives/16-asians-lookbacks-cliquets-and-forward-starts/lookback.csv};
\addplot[omDef,thick,dotted] table[x=n,y=vanilla]{figdata/derivatives/16-asians-lookbacks-cliquets-and-forward-starts/lookback.csv};
\legend{Monte Carlo,extrême décalé,continu,call vanille}
\end{axis}
\end{tikzpicture}

\omcaption{Le call lookback à strike flottant (un an, 20~\%, $r=3\%$, $q=1\%$) selon la fréquence de suivi~: simulation, la
formule continue de Goldman, Sosin et Gatto, et la même avec l'extrême décalé, traité comme une barrière. Données~: le
tutoriel.}\label{fig:dv:asians-lookbacks-cliquets-and-forward-starts:lookback}
\end{omfigure}

La couverture d'un lookback est une échelle d'options à barrière. Chaque fois que le cours atteint un nouveau plus bas, le strike est réinitialisé et la position du desk
dans l'option du dessous change. Le payoff est continu, donc il n'y a pas de saut à un niveau unique, mais le delta
saute à chaque fois qu'un nouvel extrême est établi, et le \omterm{def:dv:greeks-and-the-hedging-pnl:greeks}{vega} est élevé. C'est davantage un produit de volatilité qu'un produit directionnel.

\section{Options à départ différé et smile forward}

\begin{definition}[Option à départ différé]\label{def:dv:asians-lookbacks-cliquets-and-forward-starts:fwdstart}
Une \emph{option à départ différé}\index{option à départ différé} est une option dont le strike est fixé à une date future $t_1$ comme
proportion $k$ du spot à cette date, $K=kS_{t_1}$, et qui expire à $t_2>t_1$. Elle paie
$(S_{t_2}-kS_{t_1})^+$ pour un call.
\end{definition}

Dans le \omterm{def:dv:black-scholes-three-ways:model}{modèle de Black--Scholes}, le prix découle de l'homogénéité. En $t_1$, l'option est un call de strike $kS_{t_1}$, valant
$S_{t_1}C_{\mathrm{BS}}(1,k,t_2-t_1)$, et l'espérance actualisée de $S_{t_1}$ est $S_0e^{-qt_1}$, donc le prix est
$S_0e^{-qt_1}C_{\mathrm{BS}}(1,k,t_2-t_1)$. Il ne dépend que de la volatilité entre $t_1$ et $t_2$ et du
smile qui prévaudra en $t_1$. Dans tout autre modèle, les volatilités implicites des \omterm{def:dv:asians-lookbacks-cliquets-and-forward-starts:fwdstart}{options à départ différé} en fonction de $k$ forment
le \omterm{def:dv:local-volatility:forward}{smile forward} du chapitre 9. Ce smile n'est pas observé aujourd'hui~: c'est la prévision du modèle.

\begin{example}[Deux smiles forward à partir d'une même surface]\label{ex:dv:asians-lookbacks-cliquets-and-forward-starts:fwdsmile}
Le \omterm{def:dv:local-volatility:model}{modèle de volatilité locale} du chapitre 9 et le \omterm{def:dv:stochastic-volatility:heston}{modèle de Heston} du chapitre 10 sont tous deux calibrés sur la surface du chapitre 9~: ils recalculent
leur call à la monnaie à un an, à 19.19~\% et 18.99~\% de volatilité, contre 19.19~\% sur le marché. Leurs smiles à un mois démarrant dans six
mois sont très différents (voir la \cref{fig:dv:asians-lookbacks-cliquets-and-forward-starts:fwdsmile}). La \omterm{def:dv:local-volatility:model}{volatilité locale} donne
23.9~\% à 94~\%, 21.5~\% à la monnaie et 21.3~\% à 106~\%, presque plat. Heston donne 24.0~\%, 18.5~\% et 20.0~\%, avec un skew
et un smile. Le smile à un mois actuel passe de 19.6~\% à 14.9~\% puis à 12.0~\%. La \omterm{def:dv:local-volatility:model}{volatilité locale} aplatit le smile futur, comme
le chapitre 9 l'a établi. Heston conserve sa forme, et Bergomi a observé que ses \omterm{def:dv:local-volatility:forward}{smiles forward} sont plus convexes que le smile actuel.
\end{example}

\begin{omfigure}
\begin{tikzpicture}
\begin{axis}[width=10.5cm,height=5cm,
  xlabel={strike (\% du spot au départ)},ylabel={volatilité implicite à 1 mois (\%)},
  tick label style={font=\small},label style={font=\small},
  xmin=93,xmax=107,ymin=10,ymax=26,grid=major,grid style={gray!25},
  legend columns=3,legend cell align=left,legend style={font=\footnotesize,at={(0.5,-0.3)},anchor=north,draw=none,fill=none,/tikz/every even column/.append style={column sep=8pt}},
  table/col sep=comma]
\addplot[gray,thick,dashed,mark=o,mark size=1.6pt] table[x=k,y=today]{figdata/derivatives/16-asians-lookbacks-cliquets-and-forward-starts/forward_smiles.csv};
\addplot[omDef,very thick,mark=square*,mark size=1.6pt] table[x=k,y=lv]{figdata/derivatives/16-asians-lookbacks-cliquets-and-forward-starts/forward_smiles.csv};
\addplot[omThm,very thick,mark=*,mark size=1.6pt] table[x=k,y=heston]{figdata/derivatives/16-asians-lookbacks-cliquets-and-forward-starts/forward_smiles.csv};
\legend{smile à 1 mois aujourd'hui,volatilité locale~: dans 6 mois,Heston~: dans 6 mois}
\end{axis}
\end{tikzpicture}

\omcaption{Le smile à un mois dans six mois, tel que prévu par deux modèles calibrés sur la même surface aujourd'hui, face au
smile à un mois d'aujourd'hui. Le \omterm{def:dv:local-volatility:forward}{smile forward} de la \omterm{def:dv:local-volatility:model}{volatilité locale} est presque plat~; celui de Heston conserve un skew et un smile. Données~: le
tutoriel.}\label{fig:dv:asians-lookbacks-cliquets-and-forward-starts:fwdsmile}
\end{omfigure}

Deux effets apparaissent dans l'exemple. Le niveau est la volatilité forward~: la structure par terme de la surface est croissante, donc la
volatilité de six à sept mois est supérieure à celle à un mois d'aujourd'hui dans les deux modèles. La forme tient à la dynamique du modèle.
En \omterm{def:dv:local-volatility:model}{volatilité locale}, le skew est une fonction du niveau du spot et du temps, fixée à la calibration, et il décroît
avec la date forward. Dans Heston, le skew provient de la corrélation du spot avec une variance encore aléatoire,
et il est régénéré à chaque date. Le \omterm{def:dv:local-volatility:forward}{smile forward} du marché ne peut être lu que dans les instruments dont le paiement en dépend~:
\omterm{def:dv:asians-lookbacks-cliquets-and-forward-starts:fwdstart}{options à départ différé}, \omterm{def:dv:asians-lookbacks-cliquets-and-forward-starts:cliquet}{cliquets}, et le skew forward implicite dans les options sur l'indice de volatilité.

\section{Cliquets}

\begin{definition}[Cliquet]\label{def:dv:asians-lookbacks-cliquets-and-forward-starts:cliquet}
Un \emph{cliquet}\index{cliquet} (ratchet) est une série d'\omterm{def:dv:asians-lookbacks-cliquets-and-forward-starts:fwdstart}{options à départ différé} sur des périodes consécutives. Dans sa forme structurée
courante, il verse la somme des rendements de période $r_i=S_{t_i}/S_{t_{i-1}}-1$, chacun étant borné entre un floor local
et un cap local, la somme étant bornée entre un floor global et un cap global.
\end{definition}

\begin{definition}[Reverse cliquet]\label{def:dv:asians-lookbacks-cliquets-and-forward-starts:reverse}
Un \emph{reverse cliquet}\index{reverse cliquet} verse un gros coupon érodé par chaque rendement de période négatif~:
$\max\bigl(0,\,C+\sum_i\min(r_i,0)\bigr)$. L'investisseur est vendeur d'une série de puts à départ différé, dont le total est plafonné.
\end{definition}

Un rendement de période borné localement par un floor et un cap, $\min(\max(r_i,f),c)$, est un call spread à départ différé moins un put
spread à départ différé. Lorsque le cap et le floor sont proches de zéro, il se comporte comme un digital à départ différé, et le chapitre 15 a montré que
le prix d'un digital est fixé par le skew. Un \omterm{def:dv:asians-lookbacks-cliquets-and-forward-starts:cliquet}{cliquet} est donc un panier de paris sur le skew forward, c'est-à-dire précisément la
grandeur sur laquelle les deux modèles de la section précédente divergent. Le \omterm{def:dv:asians-lookbacks-cliquets-and-forward-starts:reverse}{reverse cliquet} est vendeur de puts à départ différé, un pari sur le
niveau de la volatilité forward et sur l'épaisseur de sa queue gauche.

\begin{example}[Le cliquet qui s'ajuste à toutes les vanilles]\label{ex:dv:asians-lookbacks-cliquets-and-forward-starts:cliquet}
Un \omterm{def:dv:asians-lookbacks-cliquets-and-forward-starts:cliquet}{cliquet} mensuel d'un an à périodes mensuelles verse la somme des douze rendements mensuels, chacun borné dans $[-1\%,+1\%]$, avec un floor à zéro au total.
Sur 100\,000 trajectoires avec des taux nuls, il vaut 1.53~\% du notionnel en \omterm{def:dv:local-volatility:model}{volatilité locale}, 2.22~\%
avec Heston et 1.19~\% avec Black--Scholes, à la volatilité à la monnaie à un an. Les trois modèles reproduisent le
call à la monnaie à un an, et les deux modèles à smile reproduisent toute la surface. Le prix de Heston dépasse de 45~\% celui en \omterm{def:dv:local-volatility:model}{volatilité locale}. Cet écart
se creuse à mesure que le cap local se resserre et que le payoff ressemble davantage à un digital~: pour $\pm0.5\%$, les prix valent 0.78~\% et 1.16~\%. Un
\omterm{def:dv:asians-lookbacks-cliquets-and-forward-starts:reverse}{reverse cliquet} qui verse $\max(0,25\%+\sum\min(r_i,0))$ vaut 5.43~\% en \omterm{def:dv:local-volatility:model}{volatilité locale}, 7.56~\% avec Heston et 3.57~\%
avec Black--Scholes (voir la \cref{fig:dv:asians-lookbacks-cliquets-and-forward-starts:cliquet}).
\end{example}

\begin{omfigure}
\begin{tikzpicture}
\begin{axis}[width=10.5cm,height=5cm,
  xlabel={cap et floor locaux ($\pm$, \%)},ylabel={prix du cliquet (\% du notionnel)},
  tick label style={font=\small},label style={font=\small},
  xmin=0,xmax=5.2,ymin=0,ymax=6.5,grid=major,grid style={gray!25},
  legend columns=3,legend style={font=\footnotesize,at={(0.5,-0.3)},anchor=north,draw=none,fill=none,/tikz/every even column/.append style={column sep=8pt}},
  table/col sep=comma]
\addplot[omThm,very thick,mark=*,mark size=1.6pt] table[x=cap,y=heston]{figdata/derivatives/16-asians-lookbacks-cliquets-and-forward-starts/cliquet_caps.csv};
\addplot[omDef,very thick,mark=square*,mark size=1.6pt] table[x=cap,y=lv]{figdata/derivatives/16-asians-lookbacks-cliquets-and-forward-starts/cliquet_caps.csv};
\addplot[gray,thick,dashed,mark=o,mark size=1.6pt] table[x=cap,y=bs]{figdata/derivatives/16-asians-lookbacks-cliquets-and-forward-starts/cliquet_caps.csv};
\legend{Heston,volatilité locale,Black--Scholes}
\end{axis}
\end{tikzpicture}

\omcaption{Un \omterm{def:dv:asians-lookbacks-cliquets-and-forward-starts:cliquet}{cliquet} mensuel d'un an (cap et floor locaux $\pm c$, floor global à zéro) sous trois modèles qui calent le
call à la monnaie à un an, et deux qui calent toute la surface. Plus le cap local est étroit, plus le payoff est une
suite de digitals à départ différé, et plus c'est le skew forward, et non le smile d'aujourd'hui, qui fixe le prix. Données~: le
tutoriel.}\label{fig:dv:asians-lookbacks-cliquets-and-forward-starts:cliquet}
\end{omfigure}

En 2004, Bergomi comptait les \omterm{def:dv:asians-lookbacks-cliquets-and-forward-starts:reverse}{reverse cliquets} et les Napoleons parmi les structures exotiques les plus récentes, et citait une revue que
Risk avait publiée en février de cette même année, sous le titre «~\omterm{def:dv:asians-lookbacks-cliquets-and-forward-starts:reverse}{Reverse cliquets}~: end of the road~?~». Son argument
est que de tels produits sont évalués par la dynamique qu'un modèle suppose pour les smiles futurs, et non par la surface d'aujourd'hui.

\section{Quel modèle pour quel payoff}

La question qu'un structurateur pose à un modèle n'est pas de savoir s'il ajuste la surface. La \omterm{def:dv:local-volatility:model}{volatilité locale}, Heston avec sauts, SABR
par échéance et les \omterm{def:dv:rough-volatility-and-forward-variance-models:bergomi}{modèles de Bergomi} l'ajustent tous, ou presque. La question est de savoir de quelles quantités au-delà de la surface le payoff
dépend, et si les hypothèses du modèle à leur sujet sont crédibles.

\begin{center}
\begin{tabular}{p{3.6cm}p{5.2cm}p{4.2cm}}
\toprule
payoff & dépend de & modèles adéquats\\
\midrule
Européenne, asiatique & distributions terminales aux dates de fixing & tout modèle qui ajuste la surface\\
à barrière, lookback & les extrêmes de la trajectoire~: le smile et sa dynamique près de la barrière & \omterm{def:dv:local-volatility:model}{volatilité locale} ou localement stochastique\\
à départ différé, \omterm{def:dv:asians-lookbacks-cliquets-and-forward-starts:cliquet}{cliquet} & le \omterm{def:dv:local-volatility:forward}{smile forward}, le skew forward & volatilité stochastique, Bergomi\\
\omterm{def:dv:asians-lookbacks-cliquets-and-forward-starts:reverse}{reverse cliquet}, Napoleon & volatilité forward et sa volatilité & Bergomi, modèles rugueux\\
options sur la variance ou sur l'indice & la volatilité de la variance & \omterm{def:dv:rough-volatility-and-forward-variance-models:fvm}{modèles de variance forward} (chapitre 12)\\
\bottomrule
\end{tabular}
\end{center}

L'exemple de Bergomi de 2004 est l'énoncé le plus clair du principe. Un Napoleon à six ans versait des coupons fixes pendant deux ans,
puis 8~\% par an augmentés de la pire des douze performances mensuelles de l'Eurostoxx 50 de chaque année, avec un plancher
à zéro. Un tel produit revient à un put sur la volatilité forward de long terme, dont un \omterm{def:dv:black-scholes-three-ways:model}{modèle de Black--Scholes} donne la valeur
sans aucune charge pour la volatilité de cette volatilité. Son modèle «~Smile dynamics II~» a été conçu pour valoriser ces produits,
le Napoleon et le \omterm{def:dv:asians-lookbacks-cliquets-and-forward-starts:reverse}{reverse cliquet} figurant parmi ses exemples, avec un contrôle séparé du skew forward court, de la
\omterm{def:dv:stochastic-volatility:sv}{corrélation spot-volatilité} et de la structure par terme de la \omterm{def:dv:stochastic-volatility:sv}{volatilité de la volatilité}. Quand deux modèles calibrés divergent de
45~\%, comme dans l'\cref{ex:dv:asians-lookbacks-cliquets-and-forward-starts:cliquet}, le desk valorise avec le modèle dont
il peut défendre la dynamique et constitue une réserve pour l'écart (chapitre 27).

\section{Tutoriel~: payoffs sur une interface commune de trajectoires}\label{tut:dv:asians-lookbacks-cliquets-and-forward-starts:tut}

\begin{tutorial}
\textbf{Objectif.} Valoriser une asiatique arithmétique avec une variable de contrôle géométrique, un lookback à observation discrète, et une
\omterm{def:dv:asians-lookbacks-cliquets-and-forward-starts:fwdstart}{option à départ différé} et un \omterm{def:dv:asians-lookbacks-cliquets-and-forward-starts:cliquet}{cliquet} sous trois modèles qui ajustent la même volatilité à la monnaie. \textbf{État final~:} les quatre
figures et les chiffres du problème du week-end.
\begin{tutsteps}
\item \textbf{La variable de contrôle.} Les deux payoffs sur les mêmes trajectoires, $\beta$ par régression, l'erreur géométrique retirée~:
\omcode{code/firm/pathdep/firm_pathdep.py}{54}{70}{Asiatique arithmétique avec la variable de contrôle géométrique.}\label{lst:dv:asians-lookbacks-cliquets-and-forward-starts:cv}
\item \textbf{Les payoffs ne lisent que des trajectoires.} Rendements périodiques, le \omterm{def:dv:asians-lookbacks-cliquets-and-forward-starts:cliquet}{cliquet} et le \omterm{def:dv:asians-lookbacks-cliquets-and-forward-starts:reverse}{reverse cliquet}~:
\omcode{code/firm/pathdep/firm_pathdep.py}{107}{121}{Payoffs de cliquet sur le tableau de trajectoires commun.}\label{lst:dv:asians-lookbacks-cliquets-and-forward-starts:cliquet}
\item \textbf{Trajectoires de chaque modèle}~: \texttt{firm\_localvol.simulate} avec les dates enregistrées, \texttt{firm\_heston.simulate\_paths},
\texttt{gbm\_paths}~; puis lancer \texttt{dv\_pathdep.forward\_smiles()}, \texttt{cliquets()} et \texttt{fig\_pathdep.py}.
\end{tutsteps}
\textbf{Pour aller plus loin.} Valoriser le \omterm{def:dv:asians-lookbacks-cliquets-and-forward-starts:cliquet}{cliquet} avec la \omterm{def:dv:stochastic-volatility:sv}{volatilité de la volatilité} de Heston divisée par deux et recalibrer $\rho$ pour conserver le
skew à un an~; ajouter un cap global de 5~\% et voir quel modèle fait le plus bouger le prix~; valoriser un coupon de type Napoleon, 8~\% plus la
pire performance mensuelle, avec un plancher à zéro.
\end{tutorial}

\section{Construction~: payoffs dépendant de la trajectoire}\label{bld:dv:asians-lookbacks-cliquets-and-forward-starts:pathdep}

\begin{build}
\bfield{Rôle} La bibliothèque de payoffs dépendant de la trajectoire de la firme miniature, écrite sur un seul tableau de trajectoires afin que tout
simulateur (Black--Scholes, \omterm{def:dv:local-volatility:model}{volatilité locale}, Heston, et le moteur du chapitre 23) valorise chacun d'eux.
\bfield{Interface} \texttt{gbm\_paths(s, times, r, q, vol, n\_paths, seed)}~; \texttt{asian\_payoff(paths, strike, right,
geometric)}, \texttt{geometric\_asian}, \texttt{asian\_mc}~; \texttt{lookback\_payoff}, \texttt{lookback\_floating\_call},
\texttt{lookback\_shifted}~; \texttt{forward\_start\_call}~; \texttt{period\_returns}, \texttt{cliquet\_payoff(paths,
local\_floor, local\_cap, global\_floor, global\_cap)}, \texttt{reverse\_cliquet\_payoff}~; \texttt{forward\_smile\_from\_paths}.
\bfield{Règles} La colonne 0 d'un tableau de trajectoires est le spot du jour~; les payoffs ne simulent jamais~; tout prix Monte-Carlo est donné avec son
erreur type~; les variables de contrôle sont utilisées dès qu'un cousin à forme fermée existe.
\bfield{Tests d'acceptation} \texttt{code/firm/pathdep/tests/}~: l'asiatique géométrique à un seul fixing coïncide avec l'option vanille et la formule
concorde avec une simulation~; la variable de contrôle divise l'erreur par dix au moins~; le lookback dépasse l'option vanille et sa
formule décalée concorde avec une simulation quotidienne~; la formule de l'\omterm{def:dv:asians-lookbacks-cliquets-and-forward-starts:fwdstart}{option à départ différé} et le \omterm{def:dv:local-volatility:forward}{smile forward} d'un modèle plat~; les payoffs de \omterm{def:dv:asians-lookbacks-cliquets-and-forward-starts:cliquet}{cliquet} sur des trajectoires
disponibles.
\bfield{Pour aller plus loin} Approximations par appariement des moments pour les asiatiques arithmétiques~; lookbacks à strike fixe~; le Napoleon.
\end{build}

\begin{omsources}
\item A.~G.~Z.~Kemna et A.~C.~F.~Vorst, «~A pricing method for options based on average asset values~», \emph{Journal of
Banking and Finance} 14(1) (1990) 113--129.
\item M.~B.~Goldman, H.~B.~Sosin et M.~A.~Gatto, «~Path dependent options~: buy at the low, sell at the high~»,
\emph{Journal of Finance} 34(5) (1979) 1111--1127.
\item L.~Bergomi, «~Smile dynamics~», \emph{Risk} (septembre 2004)~; «~Smile dynamics II~», SSRN 1493302.
\item C.~Jeffery, «~Reverse cliquets~: end of the road~?~», \emph{Risk} (février 2004) 20--22.
\end{omsources}

\section{Exercices}

\begin{exercise}[$\star$]\label{exo:dv:asians-lookbacks-cliquets-and-forward-starts:1}
Calculer le facteur de volatilité $\sqrt{(n+1)(2n+1)/(6n^2)}$ pour $n=4$ et sa limite lorsque $n\to\infty$.
\end{exercise}

\begin{exercise}[$\star$]\label{exo:dv:asians-lookbacks-cliquets-and-forward-starts:2}
Pourquoi l'asiatique arithmétique vaut-elle plus que l'asiatique géométrique~?
\end{exercise}

\begin{exercise}[$\star$]\label{exo:dv:asians-lookbacks-cliquets-and-forward-starts:3}
Sous le \omterm{def:dv:black-scholes-three-ways:model}{modèle de Black--Scholes} avec $q=2\%$, que vaut un call à départ différé six mois dans six mois, à la monnaie, par rapport à un
call à six mois à la monnaie~?
\end{exercise}

\begin{exercise}[$\star\star$]\label{exo:dv:asians-lookbacks-cliquets-and-forward-starts:4}
Montrer que $\min(\max(r,f),c)=f+(r-f)^+-(r-c)^+$ et expliquer pourquoi un \omterm{def:dv:asians-lookbacks-cliquets-and-forward-starts:cliquet}{cliquet} à cap et floor locaux étroits se comporte comme une
suite de digitals à départ différé.
\end{exercise}

\begin{exercise}[$\star\star$]\label{exo:dv:asians-lookbacks-cliquets-and-forward-starts:5}
Expliquer pourquoi le \omterm{def:dv:local-volatility:forward}{smile forward} en \omterm{def:dv:local-volatility:model}{volatilité locale} est plus plat que celui de Heston dans
l'\cref{ex:dv:asians-lookbacks-cliquets-and-forward-starts:fwdsmile}.
\end{exercise}

\begin{exercise}[$\star\star$]\label{exo:dv:asians-lookbacks-cliquets-and-forward-starts:6}
Que signifie un facteur de réduction de variance de 1\,357, en nombre de trajectoires et en temps de calcul, pour un prix souhaité à 0.001~?
\end{exercise}

\begin{exercise}[$\star\star\star$]\label{exo:dv:asians-lookbacks-cliquets-and-forward-starts:7}
\emph{Programmation.} Valoriser le \omterm{def:dv:asians-lookbacks-cliquets-and-forward-starts:reverse}{reverse cliquet} de l'\cref{ex:dv:asians-lookbacks-cliquets-and-forward-starts:cliquet} avec des coupons de
15~\% et 35~\% sous les trois modèles. Quel modèle a le prix le plus sensible au coupon~?
\end{exercise}

\begin{exercise}[$\star\star\star$]\label{exo:dv:asians-lookbacks-cliquets-and-forward-starts:8}
\emph{Trouver l'erreur.} «~Notre modèle à \omterm{def:dv:local-volatility:model}{volatilité locale} recalcule chaque vanille de la surface à un dixième de point près, donc ses prix de
\omterm{def:dv:asians-lookbacks-cliquets-and-forward-starts:cliquet}{cliquet} sont des prix de marché.~»
\end{exercise}

\section{Problème~: Le cliquet compatible avec chaque vanille}

\begin{problem}[{Problème du week-end --- une surface, deux prix}]\label{pb:dv:asians-lookbacks-cliquets-and-forward-starts:1}
Un structureur doit coter un \omterm{def:dv:asians-lookbacks-cliquets-and-forward-starts:cliquet}{cliquet} mensuel d'un an (cap et floor locaux à $\pm1\%$, floor global à zéro) sur l'indice de
la surface du chapitre 9. Le desk dispose de deux modèles calibrés~: la \omterm{def:dv:local-volatility:model}{volatilité locale} du chapitre 9 et le \omterm{def:dv:stochastic-volatility:heston}{modèle de Heston} du chapitre 10.

\medskip\noindent\textbf{Partie I --- Échauffement~: moyennes et extrêmes.}
\begin{enumerate}
\item Donner le call asiatique à douze constatations sur un an et la vanille (spot 100, 20~\%, $r=3\%$, $q=1\%$).
\item Donner le facteur de variance et la volatilité de la moyenne géométrique.
\item Donner les erreurs types avec et sans variable de contrôle.
\item Donner l'\omterm{def:dv:asians-lookbacks-cliquets-and-forward-starts:lookback}{option lookback} d'achat à strike flottant en continu et sa valeur avec surveillance quotidienne.
\item Pourquoi le \omterm{def:dv:barriers-and-digitals:shift}{décalage de barrière} échoue-t-il avec quatre dates de surveillance~?
\end{enumerate}

\medskip\noindent\textbf{Partie II --- \omterm{def:dv:local-volatility:forward}{Smiles forward}.}
\begin{enumerate}[resume]
\item Vérifier la volatilité à la monnaie à un an des deux modèles par rapport à celle du marché.
\item Donner, pour chaque modèle, la volatilité à la monnaie à un mois démarrant dans six mois, et celle d'aujourd'hui à un mois.
\item Donner, pour chaque modèle, le skew forward entre 94~\% et 106~\%, et celui d'aujourd'hui.
\item Pourquoi la volatilité forward à la monnaie est-elle supérieure à la volatilité d'aujourd'hui à un mois dans les deux modèles~?
\item Lequel des deux modèles devrait reproduire le \omterm{def:dv:local-volatility:forward}{smile forward} d'un marché d'\omterm{def:dv:asians-lookbacks-cliquets-and-forward-starts:fwdstart}{options à départ différé}, et pourquoi~?
\end{enumerate}

\medskip\noindent\textbf{Partie III --- Le \omterm{def:dv:asians-lookbacks-cliquets-and-forward-starts:cliquet}{cliquet}.}
\begin{enumerate}[resume]
\item Donner le prix du \omterm{def:dv:asians-lookbacks-cliquets-and-forward-starts:cliquet}{cliquet} sous la \omterm{def:dv:local-volatility:model}{volatilité locale}, sous Heston et sous Black--Scholes.
\item Donner les prix avec un cap et un floor locaux de $\pm0.5\%$ et $\pm5\%$.
\item Pourquoi l'écart se réduit-il lorsque la largeur du cap augmente~?
\item Donner les prix du \omterm{def:dv:asians-lookbacks-cliquets-and-forward-starts:reverse}{reverse cliquet} avec un coupon de 25~\%.
\item Quel modèle préférerait un vendeur de \omterm{def:dv:asians-lookbacks-cliquets-and-forward-starts:reverse}{reverse cliquets}, et pourquoi cette préférence constitue-t-elle un signal d'alerte~?
\end{enumerate}

\medskip\noindent\textbf{Partie IV --- Jugement.}
\begin{enumerate}[resume]
\item Quels instruments de marché permettraient de savoir quel modèle est le bon~?
\item Comment couvrir l'exposition au skew forward du \omterm{def:dv:asians-lookbacks-cliquets-and-forward-starts:cliquet}{cliquet}~?
\item Quelle réserve faudrait-il constituer si l'on valorisait avec la \omterm{def:dv:local-volatility:model}{volatilité locale}~?
\item Énoncer le \emph{résultat nommé}~: l'écart entre le prix en \omterm{def:dv:local-volatility:model}{volatilité locale} et le prix de Heston du \omterm{def:dv:asians-lookbacks-cliquets-and-forward-starts:cliquet}{cliquet} mensuel
à un an lorsque les deux modèles sont calibrés sur la même surface.
\item En une phrase~: que valorise un \omterm{def:dv:asians-lookbacks-cliquets-and-forward-starts:cliquet}{cliquet} que la vanille ne valorise pas~?
\end{enumerate}
\end{problem}

\section{Questions d'entretien}

\begin{interviewq}[$\star$ \iqroles{trader, chercheur}]\label{iq:dv:asians-lookbacks-cliquets-and-forward-starts:1}
Pourquoi une \omterm{def:dv:asians-lookbacks-cliquets-and-forward-starts:asian}{option asiatique} est-elle moins chère qu'une vanille, et de combien environ~?
\end{interviewq}

\begin{interviewq}[$\star$ \iqroles{développeur}]\label{iq:dv:asians-lookbacks-cliquets-and-forward-starts:2}
Comment valoriseriez-vous efficacement une \omterm{def:dv:asians-lookbacks-cliquets-and-forward-starts:asian}{option asiatique} arithmétique par Monte Carlo~?
\end{interviewq}

\begin{interviewq}[$\star\star$ \iqroles{chercheur}]\label{iq:dv:asians-lookbacks-cliquets-and-forward-starts:3}
Valorisez une \omterm{def:dv:asians-lookbacks-cliquets-and-forward-starts:fwdstart}{option à départ différé} sous Black--Scholes. De quoi dépend-elle dans un modèle à smile~?
\end{interviewq}

\begin{interviewq}[$\star\star$ \iqroles{trader}]\label{iq:dv:asians-lookbacks-cliquets-and-forward-starts:4}
À quoi un \omterm{def:dv:asians-lookbacks-cliquets-and-forward-starts:cliquet}{cliquet} est-il exposé, et pourquoi deux modèles calibrés sur la même surface peuvent-ils le valoriser différemment~?
\end{interviewq}

\begin{interviewq}[$\star\star$ \iqroles{risque, chercheur}]\label{iq:dv:asians-lookbacks-cliquets-and-forward-starts:5}
Pourquoi la \omterm{def:dv:local-volatility:model}{volatilité locale} est-elle un mauvais modèle pour les \omterm{def:dv:asians-lookbacks-cliquets-and-forward-starts:cliquet}{cliquets}~?
\end{interviewq}

\begin{interviewq}[$\star\star\star$ \iqroles{trader, risque}]\label{iq:dv:asians-lookbacks-cliquets-and-forward-starts:6}
Vous êtes vendeur d'un portefeuille de \omterm{def:dv:asians-lookbacks-cliquets-and-forward-starts:reverse}{reverse cliquets}. Quels sont vos risques, et comment les couvrez-vous~?
\end{interviewq}
